\documentclass[11pt]{article}
\usepackage{fontspec}
\usepackage[french,english]{babel}
\usepackage{amsmath}
\usepackage{amssymb}
\usepackage{amsthm}
\usepackage{mathrsfs}
\usepackage{tikz-cd}
\usepackage{enumitem}

\newcommand{\printedpage}[1]{\marginpar{\footnotesize p.~#1}}
\newcommand{\msp}[1]{\marginpar{\footnotesize [#1]}}
\newcommand{\Ob}{\operatorname{Ob}}
\newcommand{\Fl}{\operatorname{Fl}}
\newcommand{\card}{\operatorname{card}}
\newcommand{\Hom}{\operatorname{Hom}}
\newcommand{\Filt}{\operatorname{Filt}}
\newcommand{\Ind}{\operatorname{Ind}}
\newcommand{\circled}[1]{\tikz[baseline=(c.base)]\node[draw,circle,inner sep=0.8pt](c){#1};}

\begin{document}

\printedpage{64}
% head: A. Grothendieck : Dérivateurs, Ch. XVIII, Édition des Universités de Montpellier II et Paris VII

\msp{page 186}

\[ \card F(x) \longrightarrow \pi'^{\pi_0} \qquad \forall\, x \in \Ob C , \] % [sic] \longrightarrow -> \leqslant, a cardinal inequality is expected (cf. the next display)

thus if $G \in \Ob \mathcal{M}_\pi \subset \Ob \Filt^{\pi^{\pi_0}}$, whence
\[ \card G(x) \leqslant \pi^{\pi_0} \qquad \forall\, x \in \Ob C , \]
we find
\[ \card \Hom(G,F) \leqslant \underbrace{\underbrace{\card(\Ob C)}_{\leqslant \pi_1}\cdot(\pi'^{\pi_0})^{(\pi^{\pi_0})}}_{\pi'^{\pi_0\cdot\pi^{\pi_0}} = \pi'^{(\pi^{\pi_0})}} , \]
whence
\[ \card \Ob \mathcal{M}_\pi/F \leqslant \underbrace{\underbrace{\card(\Ob \mathcal{M}_\pi)}_{\leqslant 2^{(\pi^{\pi_0})}}\cdot\pi'^{(\pi^{\pi_0})}}_{\gamma = \pi'^{\pi^{\pi_0}}} , \]
and finally
\[ \card \Fl \mathcal{M}_\pi/F \leqslant \underbrace{\underbrace{\gamma\cdot\gamma}_{\gamma = \pi'^{(\pi^{\pi_0})}}\cdot\underbrace{\card(\Fl \mathcal{M}_\pi)}_{\leqslant 2^{(\pi^{\pi_0})}}}_{\pi'^{(\pi^{\pi_0})}} , \]
which is the announced upper bound (the same as the one proved in 4.5.12, the proof being essentially the same).

\msp{page 187}

\textbf{4.7.20.} A canonical fcg is called \textit{maximal} (where \textit{strictly canonical}) if it is not deduced from a strictly more extensive canonical fcg $(\mathcal{F}^\pi)_{\pi\geqslant\pi_{00}}$ (thus with $\pi_{00} < \pi_0$). So if $\mathcal{M}$ admits an fcga, then it even admits a strictly canonical fcg, and the latter is moreover unique. We will come back to this in the next section. % [sic] où strictement -> ou strictement, alternative is meant; [sic] n'est par -> n'est pas, negation

\section*{4.8\quad $\mathfrak{U}$-accessible categories}

\textbf{Proposition 4.8.1.} \textit{Let $\mathcal{M}$ be a $\mathfrak{U}$-category, $\pi_0$ an infinite cardinal. Consider the following conditions:}
\begin{itemize}[label={}]
\item[\textit{a)}] \textit{$\mathcal{M}$ satisfies $L_{\pi_0}$ (i.e. is stable under small $\varinjlim$ large relative to $\pi_0$), and contains a small full subcategory $S$} strictly generating, \textit{consisting of $\pi_0$-accessible objects.}
\item[\textit{b)}] \textit{There exists a small category $C$, and an equivalence of categories}
\[ \mathcal{M} \simeq \Ind_{\pi_0}(C) . \]
\item[\textit{c)}] \textit{$\mathcal{M}$ admits an accessible generalized cardinal filtration $(\mathcal{F}^\pi)_{\pi\geqslant\pi_0}$, having $\pi_0$ as initial cardinal, and such that $\pi_0$ is a useful cardinal for the said filtration.}
\end{itemize}
% foot: Version 17 décembre 2025

\printedpage{65}
% head: A. Grothendieck : Dérivateurs, Ch. XVIII, Édition des Universités de Montpellier II et Paris VII

\msp{page 188}
\begin{itemize}[label={}]
\item[\textit{c$'$)}] \textit{As in c), with $\pi_0$ a regular cardinal (and not only useful) for the accessible generalized cardinal filtration in question, and this filtration being moreover assumed} canonical.
\end{itemize}
\textit{We then have the diagram of implications}
\[ c' \Longleftrightarrow c \Longleftrightarrow b \Longrightarrow a . \]

\textbf{4.8.1.1.} I must here make the terminology precise, as to what I mean by a (full) subcategory $C$ \textit{strictly generating} in $\mathcal{M}$. For this, I refer to SGA 4 I 7.2, (i), to the equivalent conditions b) and c) of loc. cit.:
\begin{itemize}[label={}]
\item[\circled{b}] $\forall\, F \in \Ob \mathcal{M}$, the canonical arrow
\[ \varinjlim_{i\in C/F} x_i \longrightarrow F \]
is an isomorphism.
\item[\circled{c}] The canonical functor
\[ \begin{array}{ccc}
\mathcal{M} & \longrightarrow & C^{\wedge} \\[1ex]
F & \longmapsto & (\underbrace{x}_{\in \Ob C} \longmapsto \underbrace{\Hom_{\mathcal{M}}(x,F)}_{\in \Ob (\mathrm{Ens})})
\end{array} \]
is fully faithful.
\end{itemize}
These conditions imply:
\begin{itemize}[label={}]
\item[a)] $C$ is generating in $\mathcal{M}$ by strict epimorphisms, i.e. for every $F$ in $\mathcal{M}$, the family of arrows $x \longrightarrow F$ with source in $C$, with target $F$, is a strict epimorphic family.
\end{itemize}
And when $\mathcal{M}$ is stable under fiber products, this condition a) is equivalent to the preceding ones. In loc. cit. I assert

\msp{page 189}

this equivalence without a restrictive hypothesis on $\mathcal{M}$ – but in a marginal annotation I assert that this assertion is \textit{false}, outside the hypothesis of the existence of fiber products in $\mathcal{M}$. Yet I have just checked the proof given in loc. cit. of the implication a $\Longrightarrow$ b, and it is correct. But perhaps the annotation results from a confusion between condition a), and the following condition, weaker a priori, and which resembles it like a brother:
\begin{itemize}[label={}]
\item[a$'$)] There exists a family of arrows $u_i : x_i \longrightarrow F$, with source in $C$, and with target $F$, which is a strict epimorphic family (cf. SGA 4 I 10.3).
\end{itemize}
One feels like saying that when one \textit{enlarges} a strict epimorphic family of arrows, by adding other arrows with the same target, it must remain a strict epimorphic family (``a fortiori'') – which gives the impression that a$'$) must imply a), hence be equivalent to it. But if one tries to prove this hypothetical property % [sic] agrandisse -> agrandit, verb form
% foot: Version 17 décembre 2025
\printedpage{66}
% head: A. Grothendieck : Dérivateurs, Ch. XVIII, Édition des Universités de Montpellier II et Paris VII
\msp{page 190}
of stability for strict epimorphic families, one sees that one only gets there if one assumes that the initial family is moreover a \textit{universal epimorphic family} (if fiber products exist in $\mathcal{M}$), or (subject to the preceding proviso) if for every $x' \longrightarrow F$, the sieve of $\mathcal{M}/F'$ inverse image of the sieve of $\mathcal{M}/F$ generated by the $x_i \longrightarrow F$ is an ``epimorphic sieve''. This does indeed suggest that there must be cases where condition a$'$) is satisfied, but not condition a) (equivalent to b and to c).

But we are then confronted with the very unpleasant situation: if $C$ is a strictly generating subcategory of $\mathcal{M}$, it does not thereby follow that a full subcategory $C'$ of $\mathcal{M}$ such that $C' \supset C$ is also strictly generating\,! Yet I have always used previously, without thinking twice about it, that if $C$ is strictly generating, the same is true of $C' \supset C$. So I shall have to be very attentive now, % [sic] qui si -> que si, conjunction
\msp{page 191}
and check whether this $\pm$ tacit error renders certain proofs inoperative, or even, makes certain statements false. Also, I would need to give an example where $C$ is strictly generating, but $C' \supset C$ is not.

Finally, on reflection, these feared troubles disappear. I am going to state a result in due form:

\textbf{Proposition 4.8.2.} \textit{Let $C$, $\mathcal{M}$ be two $\mathfrak{U}$-categories, $i : C \longrightarrow \mathcal{M}$ a functor, whence a functor}
\[ \left\{\begin{aligned}
&\varphi : \mathcal{M} \longrightarrow C^{\wedge}\\
&\varphi(X) = (x \longmapsto \Hom_{\mathcal{M}}(i(x),X))
\end{aligned}\right. \tag*{(4.8.2.1)} \]
\textit{(composite functor $\mathcal{M} \xrightarrow{\textit{can.}} \mathcal{M}^{\wedge} \underbrace{\xrightarrow{i^*}}_{\substack{\textit{i.e. composition}\\ \textit{with } i}} C^{\wedge}$). For every $X$ in $\mathcal{M}$, we consider the category $C/X \simeq C/\varphi(X)$, and the canonical injection} [\,?] \textit{(in $\mathcal{M}^{\vee}$)}
\[ \varinjlim_{x\in C/X} i(x) \longrightarrow X . \tag*{(4.8.2.2)} \]
\begin{itemize}[label={}]
\item[\textit{1°)}] \textit{Consider the following conditions.}
\begin{itemize}[label={}]
\item[\circled{a}] \textit{For every $X \in \Ob \mathcal{M}$, the family of arrows $i(x) \longrightarrow X$, indexed by $\Ob C/X$, is a strict epimorphic family.}

\msp{page 192}
\item[\circled{b}] \textit{For every $X \in \Ob \mathcal{M}$, the canonical arrow 4.8.2.2. is an isomorphism, i.e. makes $X$ a $\varinjlim$ of the canonical functors $C/X \longrightarrow \mathcal{M}$.}
\item[\circled{c}] \textit{The functor $\varphi$ (4.8.2.1) is fully faithful.}
\end{itemize}
\textit{We then have the implications}
\[ \mathrm{b} \Longleftrightarrow \mathrm{c} \Longrightarrow \mathrm{a} , \]
\textit{and if we assume $i$ fully faithful, then the three conditions} a, b, c \textit{are equivalent.} (\textbf{N.B.} We say that the functor $i : C \longrightarrow \mathcal{M}$ is \textit{strictly generating}, if the equivalent conditions b, c are satisfied. If $C$ is a subcategory of $\mathcal{M}$ (not
% foot: Version 17 décembre 2025
\printedpage{67}
% head: A. Grothendieck : Dérivateurs, Ch. XVIII, Édition des Universités de Montpellier II et Paris VII
necessarily full), we say that $C$ is \textit{strictly generating} if the inclusion functor is strictly generating, i.e. satisfies the equivalent conditions b), c). For this, it is therefore necessary that condition a) be satisfied, i.e. that for every $X$ in $\mathcal{M}$, the family of arrows $x \longrightarrow X$ with source in $C$, with target $X$, be a strict epimorphic family, and this condition is also sufficient when $C$ is a full subcategory of $\mathcal{M}$.)
\item[\textit{2°)}] (\footnote[80]{\texttt{[2° is on page 193.]}}). \textit{Suppose that $i$ is strictly generating} (\footnote[81]{It is even enough that it satisfy the weaker condition a) of 1°).}). \textit{Then for every object $X$ of $\mathcal{M}$, the family of arrows $i(x) \longrightarrow X$, indexed by $\Ob C/X$, is not only a strict epimorphic family (condition a) in 1°), but moreover, every family of arrows with target $X$, including the preceding family, is also a strict epimorphic family.}
\item[\textit{3°)}] \textit{Consider a commutative diagram of functors of $\mathfrak{U}$-categories:}

\msp{page 193}
\[ \begin{tikzcd}[column sep=small]
C \arrow[rr, "\alpha"] \arrow[dr, "i"'] & & C' \arrow[dl, "i'"] \\
& E &
\end{tikzcd} \]
\textit{If the functor $i$ is strictly generating} (\footnote[82]{In fact, it is enough that it satisfy the weaker condition c) of 1°).}), \textit{and if moreover $i'$ is fully faithful, then $i'$ is strictly generating.}
\end{itemize}

\textbf{Corollary 4.8.3.} \textit{Let $\mathcal{M}$ be a $\mathfrak{U}$-category, $C$ a strictly generating subcategory, $C'$ a} full \textit{subcategory of $\mathcal{M}$ which contains $C$, then $C'$ is strictly generating.}

This is indeed a special case of 4.8.2 3°.

\msp{page 194}

\textsc{Proof of 4.8.2}. The proof of 1°) is that of SGA 4 I 7.2 (i). (Where one assumes that $C$ is a full subcategory of $\mathcal{M}$, but the proof given in fact proves the statement given here.) Let us prove 2°). Let us therefore add to the family of
\[ i(x) \longrightarrow X \qquad (x \in \Ob C/X) \]
a family
\[ X_\alpha \xrightarrow{p_\alpha} X \qquad (\alpha \in A) \]
of new arrows, and prove that the total family thus obtained is still a strict epimorphic family. We must therefore prove that if $Y$ is an object of $\mathcal{M}$, and if we have a family of arrows $u_x : i(x) \longrightarrow Y$, $u_\alpha : X_\alpha \longrightarrow Y$, satisfying the usual compatibility condition for pairs of arrows:
\[ (v_x,v_y) \qquad\qquad \begin{tikzcd}
& i(x) \arrow[dr, dotted, "p_x"] & \\
Z \arrow[ur, "v_x"] \arrow[dr, "v_y"'] & & X\ , \\
& i(y) \arrow[ur, dotted, "p_y"'] &
\end{tikzcd} \]
or
% foot: Version 17 décembre 2025
\printedpage{68}
% head: A. Grothendieck : Dérivateurs, Ch. XVIII, Édition des Universités de Montpellier II et Paris VII
\[ (v_x,v_\alpha) \qquad\qquad \begin{tikzcd}
& i(x) \arrow[dr, dotted, "p_x"] & \\
Z \arrow[ur, "v_x"] \arrow[dr, "v_\alpha"'] & & X\ , \\
& X_\alpha \arrow[ur, dotted, "p_\alpha"'] &
\end{tikzcd} \]
\msp{page 195}
or
\[ (v_\alpha,v_\beta) \qquad\qquad \begin{tikzcd}
& X_\alpha \arrow[dr, dotted] & \\
Z \arrow[ur, "v_\alpha"] \arrow[dr, "v_\beta"'] & & X\ , \\
& X_\beta \arrow[ur, dotted] &
\end{tikzcd} \]
namely the relations
\[ \left\{\begin{aligned}
u_xv_x &= u_yv_y\\
u_xv_x &= u_\alpha v_\alpha\\
u_\alpha v_\alpha &= u_\beta v_\beta ,
\end{aligned}\right. \tag*{($*$)} \]
then there exists an arrow
\[ u : X \longrightarrow Y \]
(necessarily unique, by property c) of strict epimorphicity of the initial family indexed by $C/X$) such that we have
\[ \left\{\begin{aligned}
up_x &= v_x \quad \text{for every } x \in \Ob C/X\\ % [sic] v_x -> u_x, the arrow u p_x goes to Y (cf. up_\alpha = u_\alpha below)
up_\alpha &= v_\alpha \quad \text{for every } \alpha \in A. % [sic] v_\alpha -> u_\alpha, cf. the relation up_\alpha = u_\alpha below
\end{aligned}\right. \tag*{($*$)} \]
Now by property a), there exists a unique $u : X \longrightarrow Y$, in such a way that the first of the relations ($*$) is satisfied, for every $x \in \Ob C/X$. So it is only a matter of proving that for this $u$, we have for every $\alpha$ the relation

\msp{page 196}
\[ up_\alpha = u_\alpha . \]
Now, applying property a) to $X_\alpha$ instead of $X$, we want that the family of arrows % [sic] veut -> voit, sense of the sentence
\[ i(y) \xrightarrow{q_y} X_\alpha \]
(indexed by $C/X_\alpha$) is a strict epimorphic family, a fortiori epimorphic. Thus to prove $up_\alpha = u_\alpha$, it is enough to prove that we have
\[ (up_\alpha)q_y = u_\alpha q_y \tag*{($**$)} \]
for every $y$ in $C/X_\alpha$. But this is a consequence of compatibility ($*$) (second formula), applied to the following case,
% foot: Version 17 décembre 2025

\end{document}
